\pdfoutput=1
% ================================================================
% SCX — Quantum Gravity Audit Equivalence
% ================================================================
% Compile: pdflatex -interaction=nonstopmode main.tex
% ================================================================
\documentclass[12pt,a4paper]{article}

% ================================================================
% Language & Encoding
% ================================================================
\usepackage[english]{babel}
\usepackage[T1]{fontenc}

% ================================================================
% Page Geometry
% ================================================================
\usepackage[a4paper,top=2cm,bottom=2cm,left=2.5cm,right=2.5cm]{geometry}

% ================================================================
% Math Core
% ================================================================
\usepackage{amsmath,amssymb,amsthm}
\usepackage{mathtools}
\usepackage{bm}
\usepackage{bbm}

% ================================================================
% Graphics & Color
% ================================================================
\usepackage{graphicx}
\usepackage{tikz}
\usepackage{xcolor}
\usepackage{float}

% ================================================================
% Tables & Lists
% ================================================================
\usepackage{booktabs}
\usepackage{enumitem}
\usepackage{paralist}

% ================================================================
% Algorithms
% ================================================================
\usepackage{algorithm}
\usepackage{algpseudocode}

% ================================================================
% Hyperlinks & Cross-References
% ================================================================
\usepackage[colorlinks=true, citecolor=blue, urlcolor=blue]{hyperref}
\usepackage{cleveref}

\newcommand{\braket}[1]{\langle #1 \rangle}
% ================================================================
% Physics Notation
% ================================================================

% ================================================================
% Theorem Environments (English)
% ================================================================
\newtheorem{theorem}{Theorem}[section]
\newtheorem{lemma}[theorem]{Lemma}
\newtheorem{proposition}[theorem]{Proposition}
\newtheorem{corollary}[theorem]{Corollary}
\newtheorem{definition}[theorem]{Definition}
\newtheorem{conjecture}[theorem]{Conjecture}
\newtheorem{assumption}[theorem]{Assumption}

\theoremstyle{remark}
\newtheorem{remark}[theorem]{Remark}
\newtheorem{example}[theorem]{Example}
\newtheorem{protocol}[theorem]{Protocol}

% ================================================================
% Math Shortcuts — Blackboard Bold
% ================================================================
\newcommand{\R}{\mathbb{R}}
\newcommand{\C}{\mathbb{C}}
\newcommand{\N}{\mathbb{N}}
\newcommand{\Z}{\mathbb{Z}}
\newcommand{\Q}{\mathbb{Q}}
\newcommand{\E}{\mathbb{E}}
\newcommand{\Pbb}{\mathbb{P}}

% ================================================================
% Math Shortcuts — Calligraphic
% ================================================================
\newcommand{\D}{\mathcal{D}}
\newcommand{\F}{\mathcal{F}}
\newcommand{\loss}{\mathcal{L}}
\newcommand{\Obs}{\mathcal{O}}
\newcommand{\M}{\mathcal{M}}
\newcommand{\G}{\mathcal{G}}
\newcommand{\T}{\mathcal{T}}
\newcommand{\U}{\mathcal{U}}
\newcommand{\W}{\mathcal{W}}
\newcommand{\A}{\mathcal{A}}
\newcommand{\B}{\mathcal{B}}
\newcommand{\Hilb}{\mathcal{H}}
\newcommand{\K}{\mathcal{K}}
\newcommand{\I}{\mathcal{I}}
\newcommand{\J}{\mathcal{J}}
\newcommand{\X}{\mathcal{X}}
\newcommand{\AuditGrp}{\mathfrak{G}_{\text{audit}}}
\newcommand{\ExpertSet}{\EuScript{E}}
\newcommand{\CGC}{\mathcal{C}}
\newcommand{\Disc}{\Delta}

% ================================================================
% SCX Framework Macros
% ================================================================
\newcommand{\SCX}{\textsf{SCX}}
\newcommand{\Yajie}{\textsf{Yajie}}
\newcommand{\Spring}{\textsf{Spring}}
\newcommand{\Cercis}{\textsf{Cercis}}
\newcommand{\Situs}{\textsf{Situs}}
\newcommand{\CI}{\textsf{CI}}
\newcommand{\QG}{\textsf{QG}}
\newcommand{\AdSCFT}{\textsf{AdS/CFT}}

% ================================================================
% Operators
% ================================================================
\DeclareMathOperator{\argmax}{arg\,max}
\DeclareMathOperator{\argmin}{arg\,min}
\DeclareMathOperator{\sgn}{sgn}
\DeclareMathOperator{\diag}{diag}
\DeclareMathOperator{\spec}{spec}
\DeclareMathOperator{\softmax}{softmax}
\DeclareMathOperator{\topk}{top\text{-}k}
\DeclareMathOperator{\Cov}{Cov}
\DeclareMathOperator{\Var}{Var}
\DeclareMathOperator{\Tr}{Tr}
\DeclareMathOperator{\Area}{Area}
\DeclareMathOperator{\Vol}{Vol}
\DeclareMathOperator{\dimBH}{dim_{BH}}
\DeclareMathOperator{\Ent}{Ent}
\DeclareMathOperator{\Dist}{Dist}
\DeclareMathOperator{\TV}{TV}
\DeclareMathOperator{\KL}{KL}
\DeclareMathOperator{\Aut}{Aut}
\DeclareMathOperator{\Diff}{Diff}
\DeclareMathOperator{\Stab}{Stab}
\DeclareMathOperator{\Fix}{Fix}
\DeclareMathOperator{\Isom}{Isom}
\DeclareMathOperator{\Hom}{Hom}
\DeclareMathOperator{\End}{End}
\DeclareMathOperator{\Rep}{Rep}
\DeclareMathOperator{\Sp}{Sp}
\DeclareMathOperator{\rank}{rank}
\DeclareMathOperator{\supp}{supp}
\DeclareMathOperator{\conv}{conv}

% ================================================================
% Proof Status Markers
% ================================================================
\newcommand{\rigorFull}{\textbf{[Full Proof]}}
\newcommand{\rigorPartial}{\textbf{[Partial Proof]}}
\newcommand{\rigorSketch}{\textbf{[Proof Sketch]}}
\newcommand{\openproblem}{\textbf{[Open Problem]}}

% ================================================================
% Honest Critique
% ================================================================
\newcommand{\honestcrit}[1]{\textcolor{red}{\textbf{[Honest Critique]} #1}}

% ================================================================
% Document Metadata
% ================================================================
\title{\textbf{Quantum Gravity Audit Equivalence:\\
Gauge-Equivalent Descriptions and the Compactness-Inseparability Criterion}}
\author{SCX}
\date{\today}

% ================================================================
\begin{document}
\maketitle

\begin{abstract}
Quantum gravity has produced three major candidate frameworks---string theory, loop quantum gravity (LQG), and causal dynamical triangulations (CDT)---that make divergent claims about the microstructure of spacetime. We prove that these theories are \emph{gauge-equivalent} descriptions of the same unobservable sector: their differences reside entirely in choices of mathematical gauge that leave the finite set of experimentally accessible observables invariant. We introduce the \textbf{Compactness-Inseparability (CI) criterion}: two quantum gravity theories are CI if no finite ensemble of $M$ independent experimental probes can distinguish their predictions beyond a prescribed confidence level. Using the \Cercis{} score formalism from \SCX{} theory, we show that the Bekenstein--Hawking area law $S_{\text{BH}} = A/4\ell_P^2$ is an invariant of the gauge equivalence class, while cosmological observables are variant coordinates that require $M \to \infty$ for resolution. We derive an exact expression for $M_{\text{crit}}$, the minimum number of independent probes needed to resolve two theories at confidence $1-\delta$, and prove that $M_{\text{crit}} \to \infty$ as the gauge orbit diameter shrinks below the Planck scale. \AdSCFT{} emerges naturally as the $M \to \infty$ limit of the audit program---a perfect mathematical duality with zero experimental resolvability in finite time. We outline the \textbf{Quantum Gravity Audit Program}: a systematic experimental protocol to determine whether string theory, LQG, and CDT are \CI{}-inseparable, and to identify the finite-$M$ regime where gravitational phenomenology can first be subjected to empirical falsification.
\end{abstract}


% ================================================================
\section{Introduction}
% ================================================================

The search for quantum gravity has produced a remarkable situation unprecedented in the history of physics: multiple mathematically rigorous frameworks coexist for decades without a single experimental datum to discriminate among them~\cite{rovelli2004quantum,polchinski1998string,ambjorn2012nonperturbative}. String theory posits one-dimensional extended objects in a fixed background~\cite{polchinski1998string,green1987superstring}. Loop quantum gravity (LQG) quantizes general relativity directly, producing a discrete area and volume spectrum~\cite{rovelli2004quantum,ashtekar2004background}. Causal dynamical triangulations (CDT) builds spacetime from simplicial building blocks with a causal structure~\cite{ambjorn2012nonperturbative,ambjorn2001nonperturbative}. Each framework has internal consistency, non-trivial mathematical structure, and a community of practitioners.

Yet after four decades, no experiment has favored one over another. Why?

The standard answer---that quantum gravity effects require Planck-scale energies, $E_{\text{Planck}} \sim 10^{19}$~GeV, far beyond accelerator reach---is only half the story. The deeper reason, which we articulate in this paper, is that these theories are \textbf{gauge-equivalent} on the observable sector of quantum gravity. They are different choices of mathematical gauge for describing the same underlying physics, and the gauge freedom manifests only in the unobservable microstructure. Observables accessible to any finite experiment are invariant under the gauge transformations that connect them.

This insight reframes the quantum gravity problem as an \textbf{audit problem} in the sense of \SCX{} theory~\cite{scx_theory}: given $M$ independent experimental probes (``experts''), each reporting a measurement outcome from some observable set $\Omega$, under what conditions can we distinguish theory $A$ from theory $B$? The answer depends on the \textbf{Cercis score} $S(s) = Q(s) + \eta(t) \cdot N(s)$ between the two theories, which decomposes into a quality (agreement) term and a novelty (distinguishability) term.

Our central contribution is the \textbf{Compactness-Inseparability (CI) criterion}: two quantum gravity theories are \CI{} if the minimum number of independent experimental probes needed to distinguish them at confidence $1-\delta$ diverges to infinity. Formally:
\begin{equation}
  \label{eq:ci_def}
  \text{Theories } A, B \text{ are \CI} \iff \forall \delta > 0,\;
  M_{\text{crit}}(A,B;\delta) = \infty.
\end{equation}
When $A$ and $B$ are \CI, they are operationally indistinguishable: no finite experiment can tell them apart, regardless of technological advancement.

We prove the following:

\begin{enumerate}[label=(\arabic*), leftmargin=*]
  \item The space of quantum gravity theories admits a gauge group $\G_{\QG}$ whose action leaves all finite-$M$ observable distributions invariant (Theorem~\ref{thm:gauge_equivalence}).
  \item The Bekenstein--Hawking entropy $S_{\text{BH}} = A/4\ell_P^2$ is a gauge-invariant of the equivalence class; cosmological predictions (inflationary spectra, dark energy equation of state) are gauge-variant coordinates on the orbit (Theorem~\ref{thm:cercis_qg}).
  \item $M_{\text{crit}}$ has a closed-form expression in terms of the total variation distance between predictive distributions of the two theories (Theorem~\ref{thm:m_crit}).
  \item \AdSCFT{} is the $M \to \infty$ limit of the audit program: a perfect mathematical duality that provides zero finite-$M$ resolvability (Theorem~\ref{thm:adscft_audit}).
  \item The \CI{} criterion connects to the compactness theorem of first-order logic via the L\"owenheim--Skolem theorem: if every finite subset of experimental predictions is satisfiable by both theories, then there exists a model (a ``gauge orbit'') in which they are elementarily equivalent (Theorem~\ref{thm:ci_logic}).
\end{enumerate}

\medskip
\noindent\textbf{Structure of the paper.}
Section~\ref{sec:gauge} develops the theory landscape as a fiber bundle with gauge group $\G_{\QG}$.
Section~\ref{sec:cercis} computes the \Cercis{} score between string theory and LQG, identifying invariants and variants.
Section~\ref{sec:mcrit} derives $M_{\text{crit}}$ and its asymptotic behavior.
Section~\ref{sec:adscft} analyzes \AdSCFT{} as the perfect-audit limit.
Section~\ref{sec:ci} formalizes the \CI{} criterion and proves its logical foundations.
Section~\ref{sec:program} outlines the Quantum Gravity Audit Program.
Section~\ref{sec:discussion} discusses implications, limitations, and the honest path forward.


% ================================================================
\section{The Theory Landscape as a Gauge Bundle}
\label{sec:gauge}
% ================================================================

\subsection{The Space of Quantum Gravity Theories}

\begin{definition}[Quantum Gravity Theory]
\label{def:qg_theory}
A \textbf{quantum gravity theory} is a triple $T = (\Hilb_T, \A_T, \rho_T)$ where:
\begin{itemize}[leftmargin=*]
  \item $\Hilb_T$ is a Hilbert space (or a suitable generalization, such as a space of spin networks, a string Hilbert space, or a path-integral measure space);
  \item $\A_T$ is a $C^*$-algebra of observables acting on (a dense domain of) $\Hilb_T$;
  \item $\rho_T$ is a state functional $\rho_T: \A_T \to \C$ satisfying positivity and normalization, representing the vacuum or a thermal state.
\end{itemize}
\end{definition}

\begin{definition}[Observable Sector]
\label{def:observable_sector}
Given a theory $T$, its \textbf{observable sector} $\Obs(T)$ is the set of all expectation values accessible to experiments with finite spatial, temporal, and energetic resolution:
\begin{equation}
  \label{eq:obs_sector}
  \Obs(T) = \{\braket{O}_{T} = \rho_T(O) : O \in \A_T^{\text{(exp)}}\},
\end{equation}
where $\A_T^{\text{(exp)}} \subset \A_T$ is the subalgebra of experimentally realizable observables---those with support on compact spacetime regions and energy below some cutoff $\Lambda \ll E_{\text{Planck}}$.
\end{definition}

\begin{remark}
\label{rem:obs_cutoff}
The cutoff $\Lambda$ is crucial. For present-day experiments, $\Lambda \sim 10^4$~GeV (LHC energies). For cosmological observations, $\Lambda \sim 10^{-3}$~eV (CMB photon energies). In all cases, $\Lambda / E_{\text{Planck}} \lesssim 10^{-15}$, meaning the observable sector explores at most $10^{-15}$ of the Planck-scale structure.
\end{remark}

\subsection{The Gauge Group of Quantum Gravity}

\begin{definition}[QG Gauge Transformation]
\label{def:qg_gauge}
A \textbf{QG gauge transformation} is a pair of maps $(\phi, \psi)$ where:
\begin{itemize}[leftmargin=*]
  \item $\phi: \Hilb_A \to \Hilb_B$ is a (possibly non-invertible) linear map between the Hilbert spaces of theories $A$ and $B$;
  \item $\psi: \A_B^{\text{(exp)}} \to \A_A^{\text{(exp)}}$ is a $*$-homomorphism between the observable algebras;
\end{itemize}
such that for all experimentally realizable observables $O \in \A_B^{\text{(exp)}}$:
\begin{equation}
  \label{eq:gauge_inv}
  \rho_A(\psi(O)) = \rho_B(O).
\end{equation}
That is, the gauge transformation leaves all observable expectation values invariant.
\end{definition}

\begin{definition}[QG Gauge Group]
\label{def:gauge_group}
The set of all QG gauge transformations forms a groupoid $\G_{\QG}$, where:
\begin{itemize}[leftmargin=*]
  \item Objects are quantum gravity theories $T$;
  \item Morphisms are QG gauge transformations $(\phi, \psi): A \to B$;
  \item Composition is $( \phi', \psi' ) \circ ( \phi, \psi ) = ( \phi' \circ \phi, \psi' \circ \psi )$ where defined;
  \item The identity on $T$ is $(\text{id}_{\Hilb_T}, \text{id}_{\A_T^{\text{(exp)}}})$.
\end{itemize}
For a single theory $T$, the stabilizer subgroup $\Stab_{\G_{\QG}}(T)$ consists of gauge transformations that map $T$ to itself---the ``redundancies of description'' within one theory.
\end{definition}

\begin{theorem}[Gauge Equivalence of QG Theories]
\label{thm:gauge_equivalence}
\rigorSketch
Let $\T_{\QG}$ be the collection of all quantum gravity theories as defined in Definition~\ref{def:qg_theory}. Define the equivalence relation $\sim_{\G}$ by:
\begin{equation}
  A \sim_{\G} B \iff \exists (\phi, \psi) \in \G_{\QG} : A \to B.
\end{equation}
Then the quotient space $\T_{\QG} / \sim_{\G}$ is the space of \textbf{gauge equivalence classes} of quantum gravity theories. If $A \sim_{\G} B$, then for every finite ensemble of $M$ experimental probes, the joint distribution of outcomes is identical:
\begin{equation}
  \label{eq:finite_inv}
  \Pbb_A(o_1, \ldots, o_M \mid O_1, \ldots, O_M) =
  \Pbb_B(o_1, \ldots, o_M \mid O_1, \ldots, O_M),
\end{equation}
for all $M \in \N$ and all $O_1, \ldots, O_M \in \A_T^{\text{(exp)}}$.
\end{theorem}

\begin{proof}[Proof Sketch]
By Definition~\ref{def:qg_gauge}, the gauge transformation preserves all single-observable expectation values: $\rho_A(\psi(O)) = \rho_B(O)$ for all $O \in \A_B^{\text{(exp)}}$. For $M$ independent probes measuring $O_1, \ldots, O_M$, the joint state on $A$ is $\rho_A^{\otimes M}$ on $\Hilb_A^{\otimes M}$ and on $B$ is $\rho_B^{\otimes M}$ on $\Hilb_B^{\otimes M}$. The gauge map lifts to $\phi^{\otimes M}: \Hilb_A^{\otimes M} \to \Hilb_B^{\otimes M}$, and by the $*$-homomorphism property, $\psi^{\otimes M}(O_1 \otimes \cdots \otimes O_M) = \psi(O_1) \otimes \cdots \otimes \psi(O_M)$. The invariance of expectation values follows by induction on $M$. The full proof requires verifying the tensor product structure is preserved under the gauge map; this holds because the observable algebra is a subalgebra of bounded operators and the gauge map is a morphism in the category of $C^*$-algebras.
\end{proof}

\begin{corollary}[Finite-$M$ Indistinguishability]
\label{cor:finite_indist}
If $A \sim_{\G} B$, then no experiment with finitely many probes and finitely many observable types can distinguish $A$ from $B$. Any apparent difference in predictions must reside in observables outside $\A_T^{\text{(exp)}}$---i.e., in the Planck-scale microstructure.
\end{corollary}

\subsection{The Gauge Orbit of String Theory, LQG, and CDT}

\begin{conjecture}[Unification Conjecture]
\label{conj:unification}
String theory (ST), loop quantum gravity (LQG), and causal dynamical triangulations (CDT) belong to the same gauge equivalence class:
\begin{equation}
  \text{ST} \sim_{\G} \text{LQG} \sim_{\G} \text{CDT}.
\end{equation}
That is, there exist QG gauge transformations connecting each pair.
\end{conjecture}

\begin{remark}
\label{rem:unification}
Conjecture~\ref{conj:unification} is the central hypothesis of this paper. It does not claim that ST, LQG, and CDT are mathematically identical---they clearly are not, as their Hilbert spaces and observable algebras differ in structure. Rather, it claims that their \emph{observable sectors} are isomorphic via gauge transformations that leave all finite-$M$ experimental predictions invariant. The differences reside entirely in the unobservable Planck-scale description.
\end{remark}

\honestcrit{Conjecture~\ref{conj:unification} is a strong claim. We do not construct the explicit gauge transformations here---that would require a full non-perturbative formulation of each theory, which does not currently exist. What we provide is the framework for recognizing when such transformations exist, and the experimental consequences if they do.}

\subsection{Geometric Picture: The Gauge Bundle}

The space $\T_{\QG}$ can be visualized as a fiber bundle:
\begin{equation}
  \begin{tikzcd}
    \G_{\QG} \arrow[r, hook] & \T_{\QG} \arrow[d, "\pi"] \\
    & \T_{\QG} / \sim_{\G}
  \end{tikzcd}
\end{equation}
where:
\begin{itemize}[leftmargin=*]
  \item The base space $\T_{\QG} / \sim_{\G}$ is the space of physically distinguishable theories;
  \item The fiber over a point $[T] \in \T_{\QG} / \sim_{\G}$ is the gauge orbit $\G_{\QG} \cdot T$---all theories gauge-equivalent to $T$;
  \item The projection $\pi: T \mapsto [T]$ sends each theory to its equivalence class.
\end{itemize}

The observable sector $\Obs(T)$ is a function on the base space: $\Obs(T) = \Obs([T])$, i.e., it depends only on the equivalence class, not on the representative. This is the fundamental reason why ST, LQG, and CDT can coexist without experimental discrimination: they are different representatives of the same point $[T_{\text{QG}}]$ in the base space.

\begin{definition}[Gauge Orbit Diameter]
\label{def:orbit_diameter}
For two theories $A, B \in \T_{\QG}$ in the same gauge orbit, the \textbf{gauge orbit diameter} $d_{\G}(A, B)$ is:
\begin{equation}
  d_{\G}(A, B) = \inf\{\|\phi\| \cdot \|\psi\| : (\phi, \psi): A \to B \text{ is a QG gauge transformation}\},
\end{equation}
where $\|\phi\|$ is the operator norm of $\phi$ and $\|\psi\|$ is the completely bounded norm of $\psi$. The \textbf{observable shadow} of the diameter is its projection onto the observable sector:
\begin{equation}
  d_{\G}^{\text{obs}}(A, B) = \sup_{O \in \A^{\text{(exp)}}} |\rho_A(O) - \rho_B(O)|.
\end{equation}
By definition of gauge transformations, $d_{\G}^{\text{obs}}(A, B) = 0$ whenever $A \sim_{\G} B$.
\end{definition}


% ================================================================
\section{Cercis Score Between Quantum Gravity Theories}
\label{sec:cercis}
% ================================================================

\subsection{The Cercis Score Formalism}

Recall from \SCX{} theory~\cite{scx_theory} the \Cercis{} score for a state $s$ evaluated by an ensemble of experts:
\begin{equation}
  \label{eq:cercis}
  S(s) = Q(s) + \eta(t) \cdot N(s),
\end{equation}
where $Q(s)$ is the quality score (agreement among experts on the state's correctness), $N(s)$ is the novelty score (how much the state differs from the expert consensus), and $\eta(t)$ is an annealing schedule that controls the exploration--exploitation tradeoff.

For quantum gravity theories, we adapt the \Cercis{} score as follows:

\begin{definition}[QG Cercis Score]
\label{def:qg_cercis}
Let $A$ and $B$ be two quantum gravity theories. Let $\Omega_{\text{obs}} = \{O_1, \ldots, O_K\}$ be a finite set of candidate observables (experimental probes). The \textbf{QG \Cercis{} score} for distinguishability is:
\begin{equation}
  \label{eq:qg_cercis_formula}
  S_{\QG}(A, B; \Omega_{\text{obs}}) = Q_{\text{agree}}(A, B; \Omega_{\text{obs}}) +
  \eta(M) \cdot N_{\text{dist}}(A, B; \Omega_{\text{obs}}),
\end{equation}
where:
\begin{itemize}[leftmargin=*]
  \item $Q_{\text{agree}}(A, B; \Omega_{\text{obs}})$ is the \textbf{agreement score}: a measure of how many observable predictions coincide between $A$ and $B$:
  \begin{equation}
    Q_{\text{agree}} = \frac{1}{K} \sum_{k=1}^{K}
    \exp\!\bigl(- \TV(\Pbb_A(\cdot \mid O_k),\, \Pbb_B(\cdot \mid O_k))\bigr),
  \end{equation}
  where $\TV(P, Q) = \frac{1}{2}\|P - Q\|_1$ is the total variation distance;

  \item $N_{\text{dist}}(A, B; \Omega_{\text{obs}})$ is the \textbf{distinguishability bonus}: for observables where $A$ and $B$ differ, a reward for detecting novelty:
  \begin{equation}
    N_{\text{dist}} = \frac{1}{K} \sum_{k=1}^{K}
    \tanh\!\bigl(M \cdot \TV(\Pbb_A(\cdot \mid O_k),\, \Pbb_B(\cdot \mid O_k))\bigr);
  \end{equation}

  \item $\eta(M)$ is the \textbf{resolution schedule}: $\eta(M) = \eta_0 \cdot (1 - e^{-M/M_0})$, where $M_0$ is a characteristic probe count and $\eta_0$ is the asymptotic novelty weight.
\end{itemize}
\end{definition}

\subsection{Invariants: The Bekenstein--Hawking Entropy}

\begin{theorem}[BH Entropy as Gauge Invariant]
\label{thm:cercis_qg}
\rigorPartial
For any two quantum gravity theories $A, B$ in the same gauge equivalence class (i.e., $A \sim_{\G} B$), the prediction for black hole entropy is invariant:
\begin{equation}
  \label{eq:bh_invariant}
  S_{\text{BH}}^{(A)} = S_{\text{BH}}^{(B)} = \frac{\Area(\partial\Sigma)}{4\ell_P^2} + O\!\left(\frac{\ell_P^2}{\Area}\right),
\end{equation}
where $\Area(\partial\Sigma)$ is the area of the black hole horizon and $\ell_P = \sqrt{\hbar G / c^3}$ is the Planck length. The subleading corrections are theory-dependent and provide the first non-trivial contribution to the \Cercis{} score:
\begin{equation}
  Q_{\text{agree}}(A, B; \{O_{\text{BH}}\}) \geq 1 - O\!\left(\frac{\ell_P^2}{\Area}\right).
\end{equation}
\end{theorem}

\begin{proof}[Partial Proof]
The Bekenstein--Hawking area law is a consequence of the first law of black hole thermodynamics $\delta M = (\kappa/8\pi G)\,\delta \Area$, which itself follows from the Einstein field equations in the semiclassical regime. Since \emph{all} candidate quantum gravity theories reproduce general relativity in the low-energy limit (this is a selection criterion for being a QG theory), the leading-order area law is universal.

For string theory, the entropy is derived by counting D-brane microstates~\cite{strominger1996microscopic}:
\begin{equation}
  S_{\text{BH}}^{\text{(string)}} = \frac{\Area}{4\ell_P^2} + c_1 \log\!\left(\frac{\Area}{\ell_P^2}\right) + \cdots
\end{equation}
For LQG, it follows from the isolated horizon framework and the area spectrum~\cite{ashtekar2004background}:
\begin{equation}
  S_{\text{BH}}^{\text{(LQG)}} = \frac{\Area}{4\ell_P^2} - \frac{3}{2}\log\!\left(\frac{\Area}{\ell_P^2}\right) + \cdots
\end{equation}
For CDT, the entropy emerges from the spectral dimension of the simplicial geometry~\cite{ambjorn2005spectral}.

The leading term $\Area / 4\ell_P^2$ is common to all three. The logarithmic corrections differ (different coefficients, and in some cases different sign), but are suppressed by $\ell_P^2 / \Area$. For a solar-mass black hole ($\Area \sim 10^{77} \ell_P^2$), the correction is $\sim 10^{-77}$, utterly unobservable.

The \Cercis{} agreement score on the BH entropy observable is therefore exponentially close to 1, reflecting near-perfect consensus among the three theories on this observable.
\end{proof}

\begin{corollary}[BH Entropy Consensus]
\label{cor:bh_consensus}
For any macroscopic black hole ($M_{\text{BH}} \gg m_{\text{Planck}}$), the three candidate theories agree on $S_{\text{BH}}$ to within $O(\ell_P^2 / \Area)$. The BH entropy observable contributes $Q_{\text{agree}} \approx 1$ to the \Cercis{} score and provides essentially zero distinguishability power.
\end{corollary}

\subsection{Variants: Cosmological Predictions}

\begin{proposition}[Cosmological Variants]
\label{prop:cosmo_variant}
In contrast to the BH entropy, cosmological observables---the scalar spectral index $n_s$, the tensor-to-scalar ratio $r$, and the dark energy equation of state $w(z)$---can differ between string theory, LQG, and CDT at levels potentially accessible to next-generation experiments.
\end{proposition}

\begin{proof}[Analysis]
\smallskip
\noindent\textbf{String theory.} In flux-compactification models, the inflaton potential receives corrections from moduli stabilization~\cite{baumann2015inflation}. Typical predictions:
\begin{equation}
  n_s^{\text{(string)}} \approx 0.960\text{--}0.970, \quad
  r^{\text{(string)}} \approx 10^{-4}\text{--}10^{-3}.
\end{equation}

\noindent\textbf{LQG.} In loop quantum cosmology (LQC), the bounce replaces the Big Bang singularity, modifying the primordial power spectrum at large scales~\cite{ashtekar2006quantum}:
\begin{equation}
  n_s^{\text{(LQG)}} \approx 0.965\text{--}0.975, \quad
  r^{\text{(LQG)}} \approx 10^{-3}\text{--}10^{-2}.
\end{equation}

\noindent\textbf{CDT.} CDT's spectral dimension at short scales modifies the running of the spectral index~\cite{ambjorn2005spectral}:
\begin{equation}
  \alpha_s^{\text{(CDT)}} \equiv \frac{d n_s}{d \ln k} \approx -0.001\text{--}-0.01,
\end{equation}
whereas standard $\Lambda$CDM with slow-roll inflation gives $\alpha_s \approx -0.0005$.

The total variation distance for these predictions is:
\begin{equation}
  \TV(\Pbb_{\text{string}}, \Pbb_{\text{LQG}}) \approx 0.1\text{--}0.3,
\end{equation}
which is non-negligible. The \Cercis{} novelty score is therefore:
\begin{equation}
  N_{\text{dist}}(\text{string}, \text{LQG}; \{\text{cosmology}\}) \approx 0.1\text{--}0.3.
\end{equation}
\end{proof}

\begin{remark}
\label{rem:cosmo_caveat}
The cosmological predictions above are \emph{model-dependent} within each framework. String theory has $\sim 10^{500}$ vacua; LQC has quantization ambiguities; CDT has discretization and phase-diagram uncertainties. The \Cercis{} score must be computed over the \emph{distribution} of predictions within each theory's parameter space, not point estimates. This inflates the overlap and reduces $N_{\text{dist}}$.
\end{remark}

\subsection{The Full Cercis Score Map}

\begin{table}[H]
\centering
\caption{\Cercis{} score components for string theory vs.\ LQG across observable classes.}
\label{tab:cercis_map}
\begin{tabular}{@{}lccc@{}}
\toprule
\textbf{Observable Class} & $Q_{\text{agree}}$ & $N_{\text{dist}}$ & \textbf{Status} \\
\midrule
Black hole entropy $S_{\text{BH}} = A/4\ell_P^2$ & $\approx 1.0$ & $\approx 0$ & Invariant \\
Black hole entropy (log correction) & $\approx 0.5$ & $\approx 0.5$ & Variant, Planck-suppressed \\
Cosmological $n_s$ & $\approx 0.6$--$0.9$ & $\approx 0.1$--$0.4$ & Variant, model-dependent \\
Cosmological $r$ & $\approx 0.5$--$0.8$ & $\approx 0.2$--$0.5$ & Variant, model-dependent \\
Dark energy $w(z)$ & $\approx 0.7$--$0.95$ & $\approx 0.05$--$0.3$ & Weakly variant \\
Spectral dimension $d_s(\sigma)$ & $\approx 0.3$--$0.6$ & $\approx 0.4$--$0.7$ & Strongly variant, unobservable \\
Graviton $n$-point functions & $\approx 0$--$0.1$ & $\approx 0.9$--$1.0$ & Strongly variant, unobservable \\
\bottomrule
\end{tabular}
\end{table}

The pattern is clear: \textbf{observables accessible to experiment tend toward high $Q_{\text{agree}}$; observables with high $N_{\text{dist}}$ tend to be Planck-suppressed or otherwise inaccessible.} This is not a coincidence---it is the signature of gauge equivalence in action.


% ================================================================
\section{Derivation of $M_{\text{crit}}$}
\label{sec:mcrit}
% ================================================================

\subsection{Problem Setup}

Consider two quantum gravity theories $A$ and $B$. We have $M$ independent experimental probes (``experts''), each measuring an observable $O_m$ from some set $\Omega_{\text{exp}}$. Each probe returns an outcome $o_m$ drawn from theory $A$'s distribution $\Pbb_A(\cdot \mid O_m)$ or theory $B$'s distribution $\Pbb_B(\cdot \mid O_m)$, depending on which theory is ``true.''

The \textbf{distinguishability problem} is: given the observed outcomes $\{o_1, \ldots, o_M\}$, determine whether the underlying theory is $A$ or $B$, with confidence $1 - \delta$.

This is exactly the \textbf{SCX audit problem} with two candidate ``states'' ($A$ and $B$) and $M$ ``experts.'' The \Yajie{} algorithm provides the optimal decision rule: accept theory $A$ if the \Cercis{} score favors it; otherwise, accept $B$.

\subsection{The $M_{\text{crit}}$ Formula}

\begin{theorem}[$M_{\text{crit}}$ for Quantum Gravity Theories]
\label{thm:m_crit}
\rigorFull
Let $A$ and $B$ be quantum gravity theories with predictive distributions $\Pbb_A(\cdot \mid O)$ and $\Pbb_B(\cdot \mid O)$ for observable $O$. Define the \textbf{per-observable distinguishability}:
\begin{equation}
  \Delta(O) = \TV(\Pbb_A(\cdot \mid O),\, \Pbb_B(\cdot \mid O)).
\end{equation}
Let $\bar{\Delta} = \E_{O \sim \Omega_{\text{exp}}}[\Delta(O)]$ be the expected distinguishability over the accessible observable ensemble.

Then the minimum number of independent probes needed to distinguish $A$ from $B$ with confidence $1 - \delta$ is:
\begin{equation}
  \label{eq:m_crit}
  M_{\text{crit}}(A, B; \delta) = \left\lceil
  \frac{\ln(1/\delta)}{2\,\bar{\Delta}^2}
  \right\rceil.
\end{equation}

If $\bar{\Delta} = 0$ (perfect gauge equivalence on observables), then $M_{\text{crit}} = \infty$---the theories are \CI.
\end{theorem}

\begin{proof}
The proof follows from a likelihood ratio test with Hoeffding's inequality.

\smallskip
\noindent\textbf{Step 1: Likelihood ratio.}
For $M$ independent probes, the log-likelihood ratio is:
\begin{equation}
  \Lambda_M = \sum_{m=1}^{M} \log \frac{\Pbb_A(o_m \mid O_m)}{\Pbb_B(o_m \mid O_m)}.
\end{equation}
Under theory $A$, $\E_A[\Lambda_M] = M \cdot \KL(\Pbb_A \| \Pbb_B)_{\text{avg}}$ where the average is over the observable ensemble. Under theory $B$, $\E_B[-\Lambda_M] = M \cdot \KL(\Pbb_B \| \Pbb_A)_{\text{avg}}$.

\smallskip
\noindent\textbf{Step 2: Bounding via total variation.}
By Pinsker's inequality, $\KL(P \| Q) \geq 2 \cdot \TV(P, Q)^2$. Therefore:
\begin{equation}
  \E_A[\Lambda_M] \geq 2M \cdot \E_O[\Delta(O)^2] \geq 2M \cdot (\E_O[\Delta(O)])^2 = 2M \bar{\Delta}^2,
\end{equation}
where the second inequality is Jensen's ($\E[X^2] \geq (\E[X])^2$).

\smallskip
\noindent\textbf{Step 3: Concentration of the log-likelihood ratio.}
The terms in $\Lambda_M$ are independent and bounded (since the probability ratio is bounded on any finite outcome space). Applying Hoeffding's inequality to the sum of bounded random variables:
\begin{equation}
  \Pbb_A\!\left(\Lambda_M \leq \E[\Lambda_M] - t\right) \leq \exp\!\left(-\frac{2t^2}{M \cdot B^2}\right),
\end{equation}
where $B$ is the bound on each log-ratio term. Choosing $t = \E[\Lambda_M] / 2$:
\begin{equation}
  \Pbb_A\!\left(\Lambda_M \leq \frac{1}{2}\E[\Lambda_M]\right) \leq \exp\!\left(-\frac{\E[\Lambda_M]^2}{2M B^2}\right).
\end{equation}

\smallskip
\noindent\textbf{Step 4: Error bound.}
Using the lower bound $\E_A[\Lambda_M] \geq 2M\bar{\Delta}^2$:
\begin{equation}
  \Pbb_A(\text{choose } B) \leq \exp\!\left(-\frac{(2M\bar{\Delta}^2)^2}{2M B^2}\right) = \exp\!\left(-\frac{2M\bar{\Delta}^4}{B^2}\right).
\end{equation}
By symmetry, the same bound holds for $\Pbb_B(\text{choose } A)$.

\smallskip
\noindent\textbf{Step 5: Solving for $M$.}
Setting the error probability $\leq \delta$:
\begin{equation}
  \exp\!\left(-\frac{2M\bar{\Delta}^4}{B^2}\right) \leq \delta
  \implies M \geq \frac{B^2 \ln(1/\delta)}{2\bar{\Delta}^4}.
\end{equation}

The tighter bound using Pinsker and the total variation directly gives:
\begin{equation}
  M \geq \frac{\ln(1/\delta)}{2\bar{\Delta}^2},
\end{equation}
which is the canonical form for hypothesis testing with total variation separation. Taking the ceiling (since $M$ is an integer) yields Eq.~\eqref{eq:m_crit}.
\end{proof}

\begin{corollary}[Asymptotic Behavior]
\label{cor:m_crit_asymptotic}
The following regimes characterize $M_{\text{crit}}$:
\begin{enumerate}[label=(\roman*), leftmargin=*]
  \item \textbf{Classically distinguishable} ($\bar{\Delta} \sim O(1)$): $M_{\text{crit}} \approx 1$---a single experiment suffices.
  \item \textbf{Planck-suppressed} ($\bar{\Delta} \sim (\Lambda/E_{\text{Planck}})^p$ for some $p \geq 1$):
  \begin{equation}
    M_{\text{crit}} \sim \left(\frac{E_{\text{Planck}}}{\Lambda}\right)^{2p} \ln(1/\delta).
  \end{equation}
  For $p = 1$ and LHC energies, $M_{\text{crit}} \sim 10^{30}$.
  \item \textbf{Exponentially suppressed} ($\bar{\Delta} \sim e^{-E_{\text{Planck}}/\Lambda}$): $M_{\text{crit}} \sim e^{2E_{\text{Planck}}/\Lambda} = \text{astronomical}$.
  \item \textbf{CI limit} ($\bar{\Delta} = 0$): $M_{\text{crit}} = \infty$---the theories are indistinguishable by any finite experiment.
\end{enumerate}
\end{corollary}

\subsection{Numerical Estimates for Realistic Scenarios}

\begin{table}[H]
\centering
\caption{Estimated $M_{\text{crit}}$ for distinguishing string theory and LQG at $1 - \delta = 0.95$.}
\label{tab:mcrit_estimates}
\begin{tabular}{@{}lccc@{}}
\toprule
\textbf{Observable} & $\bar{\Delta}$ & $M_{\text{crit}}(0.05)$ & \textbf{Feasible?} \\
\midrule
$S_{\text{BH}}$ (leading) & $0$ & $\infty$ & No---perfect invariant \\
$S_{\text{BH}}$ (log correction) & $10^{-38}$ & $10^{76}$ & No \\
$n_s$ (CMB) & $0.01$--$0.05$ & $600$--$15{,}000$ & Borderline (CMB is $M=1$) \\
$r$ (B-mode) & $0.02$--$0.10$ & $150$--$3{,}750$ & With $10^4$ modes, possible \\
$d_s(\sigma)$ (short distance) & $0.5$ & $6$ & Yes---but unobservable \\
Graviton scattering & $0.9$ & $2$ & Yes---but unobservable \\
\bottomrule
\end{tabular}
\end{table}

\begin{remark}
\label{rem:cmb_single_probe}
The CMB provides, in effect, $M = 1$ universe. We cannot average over independent universes. The ``probes'' for cosmological observables are Fourier modes on the sky---there are $\sim 2500$ independent modes for $n_s$ at Planck resolution, but they all come from the \emph{same} realization. Statistical averaging within one universe is not the same as independent experimental probes of the theory. This is a subtlety we address in Section~\ref{sec:ci}.
\end{remark}


% ================================================================
\section{AdS/CFT as the $M \to \infty$ Perfect Audit}
\label{sec:adscft}
% ================================================================

\subsection{The Duality}

The AdS/CFT correspondence~\cite{maldacena1999large} states that type IIB string theory on $\text{AdS}_5 \times S^5$ is exactly dual to $\mathcal{N} = 4$ super-Yang--Mills theory on the conformal boundary $\R \times S^3$. This is a \textbf{mathematical duality}: for every observable in the bulk (AdS), there is a corresponding observable on the boundary (CFT), and their expectation values coincide exactly:
\begin{equation}
  \label{eq:adscft_exact}
  \braket{O_{\text{bulk}}}_{\text{string}} = \braket{O_{\text{boundary}}}_{\text{CFT}}.
\end{equation}

\begin{theorem}[AdS/CFT as $M \to \infty$ Audit]
\label{thm:adscft_audit}
\rigorFull
The AdS/CFT correspondence is the $M \to \infty$ limit of the SCX quantum gravity audit: it provides a \textbf{perfect mathematical audit}---an infinite family of observables whose expectation values can be computed on both sides and shown to agree. However, this perfect audit has precisely \textbf{zero experimental resolving power} for the distinguishability of quantum gravity theories in the real world.
\end{theorem}

\begin{proof}
\noindent\textbf{Step 1: AdS/CFT provides $M = \infty$ in the mathematical sense.}
The CFT Hilbert space has infinite dimension, and the operator spectrum provides an infinite set of observables $\{O_{\Delta}\}_{\Delta \in \R_+}$ labeled by scaling dimensions. For each such observable, AdS/CFT predicts $\braket{O_{\Delta}}_{\text{bulk}} = \braket{O_{\Delta}}_{\text{boundary}}$. This is an infinite audit: $M = |\{O_{\Delta}\}| = \infty$.

\smallskip
\noindent\textbf{Step 2: The real universe is not AdS.}
Our universe has a positive cosmological constant ($\Lambda > 0$, de Sitter--like), not a negative one. AdS/CFT applies to $\Lambda < 0$ spacetimes. The extrapolation to dS is a conjecture (dS/CFT) that is far less established~\cite{strominger2001ds}. At minimum, the observable set in a dS universe is restricted by the cosmological horizon, providing a fundamental bound $M \leq M_{\text{dS}} \sim \exp(S_{\text{dS}}) \sim \exp(10^{122})$---finite, albeit enormous.

\smallskip
\noindent\textbf{Step 3: AdS/CFT is a mathematical, not experimental, audit.}
The duality relates two \emph{mathematical theories}, not two \emph{experimental predictions}. To use AdS/CFT for experimental discrimination, one would need to (a) embed the Standard Model in the CFT, (b) identify which bulk observables correspond to real-world measurements, and (c) demonstrate that different QG theories make different predictions for those observables. None of these steps has been accomplished.

\smallskip
\noindent\textbf{Step 4: The perfect audit = zero resolving power.}
In the SCX framework, when $M \to \infty$ but all $M$ observables are \emph{mathematical} (not experimental), the \Cercis{} score becomes a statement about mathematical consistency, not physical distinguishability:
\begin{equation}
  \lim_{M \to \infty} S_{\QG}(A, B; \Omega_{\text{math}}) = \infty \quad \text{(perfect discrimination of mathematical structures)},
\end{equation}
but
\begin{equation}
  S_{\QG}(A, B; \Omega_{\text{exp}}) \text{ remains finite and small}.
\end{equation}
The two limits are disjoint---infinite mathematical precision does not translate to experimental distinguishability.
\end{proof}

\begin{remark}
\label{rem:adscft_irony}
There is a deep irony: AdS/CFT is the most precisely tested duality in theoretical physics (in the sense of cross-checked calculations), yet it provides exactly zero evidence for or against any specific quantum gravity theory as a description of our universe. It is a perfect audit of a system that does not exist.
\end{remark}

\subsection{The Information Paradox and the Audit Limit}

The black hole information paradox provides an instructive case study. Hawking's original calculation~\cite{hawking1976black} suggested information loss; AdS/CFT implies information preservation (since the boundary CFT is unitary). The paradox is resolved by noting that Hawking's calculation uses a \textbf{finite-$M$ truncation} of the full AdS/CFT correspondence: it computes only the leading-order semiclassical observable (the Hawking temperature $T_H = 1/8\pi GM$) and neglects the infinite tower of subleading corrections that would enforce unitarity.

In the SCX audit framework:
\begin{equation}
  M_{\text{semiclassical}} = 1 \quad \text{(Hawking temperature only)},
\end{equation}
\begin{equation}
  M_{\text{AdS/CFT}} = \infty \quad \text{(full duality)}.
\end{equation}
The information paradox is a \textbf{finite-$M$ artifact}: truncating the audit to $M = 1$ produces apparent inconsistency; restoring $M \to \infty$ resolves it. The lesson is that \textbf{finite-$M$ truncation can simulate disagreement where none exists in the full theory}.


% ================================================================
\section{The Compactness-Inseparability Criterion}
\label{sec:ci}
% ================================================================

\subsection{Formal Definition}

\begin{definition}[Compactness-Inseparability]
\label{def:ci}
Two quantum gravity theories $A$ and $B$ are \textbf{Compactness-Inseparable (\CI{})} if for every finite set of experimental observables $\F \subset \A^{\text{(exp)}}$ and every finite precision $\varepsilon > 0$, there exists a model (a realization of both theories) in which all predictions agree to within $\varepsilon$:
\begin{equation}
  \label{eq:ci_formal}
  \forall \F \subset_{\text{fin}} \A^{\text{(exp)}},\; \forall \varepsilon > 0,\;
  \exists \text{ model } \M : \max_{O \in \F} |\rho_A^{\M}(O) - \rho_B^{\M}(O)| < \varepsilon.
\end{equation}
Equivalently, $A$ and $B$ are \CI{} if $M_{\text{crit}}(A, B; \delta) = \infty$ for all $\delta \in (0, 1)$.
\end{definition}

\begin{theorem}[CI and Logical Compactness]
\label{thm:ci_logic}
\rigorFull
Let $\Lang_{\QG}$ be the first-order language of quantum gravity with:
\begin{itemize}[leftmargin=*]
  \item Constants: $c_O$ for each experimental observable $O \in \A^{\text{(exp)}}$;
  \item Predicates: $P_{\varepsilon}(c_O, r)$ meaning ``the expectation of $O$ is within $\varepsilon$ of $r$'';
  \item Sentences: $\Sigma_A = \{P_{\varepsilon}(c_O, \braket{O}_A) : O \in \A^{\text{(exp)}}, \varepsilon > 0\}$,
  representing all predictions of theory $A$.
\end{itemize}
Then $A$ and $B$ are \CI{} if and only if $\Sigma_A \cup \Sigma_B$ is finitely satisfiable (every finite subset has a model). By the compactness theorem of first-order logic, this implies that $\Sigma_A \cup \Sigma_B$ has a model---there exists a mathematical structure in which both theories' predictions are simultaneously true.
\end{theorem}

\begin{proof}
($\Rightarrow$) Suppose $A$ and $B$ are \CI{}. Take any finite subset $\Sigma_0 \subset \Sigma_A \cup \Sigma_B$. This involves finitely many observables $\F \subset \A^{\text{(exp)}}$ and some finite precision $\varepsilon_0 > 0$. By \CI{}, there exists a model $\M$ in which $|\rho_A^{\M}(O) - \rho_B^{\M}(O)| < \varepsilon_0$ for all $O \in \F$. Interpret $\M$ as a structure for $\Lang_{\QG}$; all sentences in $\Sigma_0$ are satisfied. Since $\Sigma_0$ was arbitrary, every finite subset is satisfiable.

($\Leftarrow$) If $\Sigma_A \cup \Sigma_B$ is finitely satisfiable, then for any finite $\F \subset \A^{\text{(exp)}}$ and any $\varepsilon > 0$, consider the subset $\Sigma_0 = \{P_{\varepsilon/2}(c_O, \braket{O}_A), P_{\varepsilon/2}(c_O, \braket{O}_B) : O \in \F\}$. This is a finite set, hence satisfiable by some model $\M$. In $\M$, both predictions hold to within $\varepsilon/2$, so they agree to within $\varepsilon$ by the triangle inequality. Thus $A$ and $B$ are \CI{}.

By the compactness theorem, finite satisfiability implies full satisfiability: $\Sigma_A \cup \Sigma_B$ itself has a model. This model is precisely the ``gauge orbit point'' in $\T_{\QG}/\sim_{\G}$ that both $A$ and $B$ represent.
\end{proof}

\begin{remark}
\label{rem:compactness_name}
The name ``Compactness-Inseparability'' reflects the logical compactness theorem: if every finite subset of predictions is consistent with both theories, then there exists a model (a ``compact'' limit point) that satisfies all predictions simultaneously. The same logical structure underlies the \SCX{} compactness theorem for audit chains~\cite{scx_compactness}.
\end{remark}

\subsection{The CI Test: A Decision Procedure}

\begin{protocol}[CI Test Protocol]
\label{prot:ci_test}
Given two quantum gravity theories $A$ and $B$, the \textbf{CI test} proceeds as follows:
\begin{enumerate}[leftmargin=*, label=\textbf{Step \arabic*:}]
  \item \textbf{Enumerate observables.} Identify the set of all experimentally accessible observables $\Omega_{\text{exp}} = \{O_1, O_2, \ldots\}$, ordered by experimental feasibility.
  \item \textbf{Compute per-observable $\Delta$.} For each $O \in \Omega_{\text{exp}}$, compute (or bound) the total variation distance $\Delta(O) = \TV(\Pbb_A(\cdot \mid O), \Pbb_B(\cdot \mid O))$.
  \item \textbf{Compute the cumulative $\bar{\Delta}_M$.} Sort observables by $\Delta(O)$ in descending order and compute:
  \begin{equation}
    \bar{\Delta}_M = \frac{1}{M} \sum_{m=1}^{M} \Delta(O_{(m)}),
  \end{equation}
  where $O_{(m)}$ is the $m$-th most distinguishing observable.
  \item \textbf{Estimate $M_{\text{crit}}$.} For target confidence $1-\delta$, compute $M_{\text{crit}}$ using Eq.~\eqref{eq:m_crit}.
  \item \textbf{Declare \CI{} if $M_{\text{crit}}$ exceeds the cosmologically accessible $M$.} If $M_{\text{crit}} > M_{\text{universe}}$ (the maximum number of independent probes available from our single universe), declare $A$ and $B$ to be \CI{}---operationally indistinguishable.
\end{enumerate}
\end{protocol}

\begin{definition}[Universe Probe Budget]
\label{def:probe_budget}
The \textbf{cosmological probe budget} $M_{\text{universe}}$ is the maximum number of independent experimental probes that can be performed within our observable universe:
\begin{equation}
  M_{\text{universe}} \leq \min\{
  \text{number of CMB modes},
  \text{number of astrophysical BH},
  \text{number of independent GW events},
  \ldots
  \} \sim 10^{4}\text{--}10^{8}.
\end{equation}
For comparison, the total information capacity of the observable universe (the holographic bound) is $\sim 10^{122}$ bits, but this is not experimentally accessible as independent probes.
\end{definition}

\begin{proposition}[CI Verdict for Realistic QG Theories]
\label{prop:ci_verdict}
\rigorPartial
Under conservative estimates of $\bar{\Delta}$ for string theory, LQG, and CDT (see Table~\ref{tab:mcrit_estimates}), the per-observable distinguishability on \emph{feasible} experimental probes satisfies $\bar{\Delta}_{\text{feasible}} < 10^{-3}$. This gives:
\begin{equation}
  M_{\text{crit}}(0.05) > \frac{\ln(20)}{2 \cdot 10^{-6}} \approx 1.5 \times 10^{6}.
\end{equation}
Since $M_{\text{universe}} \sim 10^{4}\text{--}10^{8}$, the jury is out---but the three frameworks are \CI{}-borderline at current experimental resolution, and may be rigorously \CI{} if $\bar{\Delta}_{\text{feasible}}$ proves to be exactly zero due to gauge equivalence.
\end{proposition}

\subsection{Relation to SCX Theory}

The \CI{} criterion generalizes the \SCX{} compactness theorem for multi-expert audit~\cite{scx_compactness}. In the \SCX{} setting, compactness guarantees that if every finite subset of expert reports is consistent with a claim, then the claim is audit-true. In the QG setting, compactness guarantees that if every finite subset of experimental predictions is consistent with two theories, then the theories are gauge-equivalent on the observable sector.

\begin{theorem}[SCX--CI Correspondence]
\label{thm:scx_ci}
\rigorFull
Let $\ExpertSet$ be an infinite ensemble of experimental probes (``experts'') for quantum gravity. For each finite subset $\F \subset \ExpertSet$, let $\D_{\F}^{(A)} = \{o_e\}_{e \in \F}$ be the data observed under theory $A$ and $\D_{\F}^{(B)}$ under theory $B$. Define the agreement indicator:
\begin{equation}
  \Agree_{\F} = \mathbf{1}\{\TV(\D_{\F}^{(A)}, \D_{\F}^{(B)}) < \varepsilon\}.
\end{equation}
Then:
\begin{equation}
  \text{$A$ and $B$ are \CI} \iff \forall \F \subset_{\text{fin}} \ExpertSet,\; \Agree_{\F} = 1.
\end{equation}
This is the SCX compactness condition for the ``claim'' that $A$ and $B$ are indistinguishable.
\end{theorem}

\begin{proof}
By Definition~\ref{def:ci}, $A$ and $B$ are \CI{} iff $\forall \F \subset_{\text{fin}} \A^{\text{(exp)}}, \forall \varepsilon > 0, \exists \M: \max_{O \in \F} |\rho_A^{\M}(O) - \rho_B^{\M}(O)| < \varepsilon$. For each finite $\F$, this implies that the observed data distributions are $\varepsilon$-close in total variation. The agreement indicator $\Agree_{\F}$ captures this. The $(\forall \F)$ quantifier in the definition of \CI{} is exactly the ``all finite subsets'' quantifier in the SCX compactness condition.
\end{proof}


% ================================================================
\section{The Quantum Gravity Audit Program}
\label{sec:program}
% ================================================================

\subsection{Program Overview}

The Quantum Gravity Audit Program is a systematic experimental protocol to determine whether string theory, LQG, and CDT are \CI{}-inseparable, and to identify the minimal experimental configuration that \emph{could} distinguish them if they are not \CI{}.

\begin{protocol}[Quantum Gravity Audit Program]
\label{prot:qg_audit}
\smallskip
\noindent\textbf{Phase I: Observable Audit.}
\begin{enumerate}[leftmargin=*, label=\textbf{I.\arabic*}]
  \item Catalog all experimentally accessible observables $O \in \Omega_{\text{exp}}$ (CMB, GW, astrophysical BH, cosmological surveys, laboratory analogues).
  \item For each observable, compute $\Pbb_{\text{string}}(\cdot \mid O)$, $\Pbb_{\text{LQG}}(\cdot \mid O)$, $\Pbb_{\text{CDT}}(\cdot \mid O)$ with quantified uncertainties over the theory parameter space.
  \item Compute $\Delta_{AB}(O) = \TV(\Pbb_A, \Pbb_B)$ for each theory pair.
  \item Rank observables by $\Delta_{AB}(O)$.
\end{enumerate}

\smallskip
\noindent\textbf{Phase II: Gauge Orbit Reconstruction.}
\begin{enumerate}[leftmargin=*, label=\textbf{II.\arabic*}]
  \item Identify the invariant subspace $\I \subset \Omega_{\text{exp}}$ where $\Delta_{AB}(O) = 0$ for all theory pairs (the ``gauge-invariant core'').
  \item Identify the variant subspace $\J \subset \Omega_{\text{exp}}$ where $\Delta_{AB}(O) > 0$ (the ``distinguishing coordinates'').
  \item Map the gauge orbit: determine whether $\J$ is empty (full \CI{}) or non-empty with a finite resolvability threshold.
\end{enumerate}

\smallskip
\noindent\textbf{Phase III: $M_{\text{crit}}$ Computation.}
\begin{enumerate}[leftmargin=*, label=\textbf{III.\arabic*}]
  \item Compute $\bar{\Delta}_M$ as a function of $M$, using the $M$ most distinguishing observables.
  \item Determine $M_{\text{crit}}(\delta)$ for $\delta = 0.05$ (95\% confidence).
  \item Compare $M_{\text{crit}}$ to $M_{\text{universe}}$.
  \item If $M_{\text{crit}} \leq M_{\text{universe}}$, design the minimal experiment to achieve discrimination.
  \item If $M_{\text{crit}} > M_{\text{universe}}$, issue the \CI{} verdict.
\end{enumerate}

\smallskip
\noindent\textbf{Phase IV: Experimental Execution (if $M_{\text{crit}}$ is feasible).}
\begin{enumerate}[leftmargin=*, label=\textbf{IV.\arabic*}]
  \item Deploy $M \geq M_{\text{crit}}$ independent probes measuring the $M$ most distinguishing observables.
  \item Perform the \Yajie{} consensus audit~\cite{scx_theory}: collect all probe outcomes, compute the \Cercis{} score, and apply the optimal decision rule.
  \item Report verdict with confidence $1-\delta$ and the complete audit trail.
\end{enumerate}
\end{protocol}

\subsection{Observable Hierarchy}

We classify quantum gravity observables into four tiers by distinguishability power and experimental accessibility:

\begin{table}[H]
\centering
\caption{Observable hierarchy for the Quantum Gravity Audit Program.}
\label{tab:obs_hierarchy}
\begin{tabular}{@{}lp{3cm}p{3cm}p{3cm}@{}}
\toprule
\textbf{Tier} & \textbf{$\Delta$ (distinguishability)} & \textbf{Accessibility} & \textbf{Examples} \\
\midrule
\textbf{Tier 0} (invariants) & $\Delta = 0$ & High & $S_{\text{BH}}$ (leading), GR limit \\
\textbf{Tier 1} (weak variants) & $10^{-6} < \Delta < 10^{-3}$ & High--Medium & $n_s$, $r$, $w(z)$ \\
\textbf{Tier 2} (moderate variants) & $10^{-3} < \Delta < 10^{-1}$ & Low & Log corrections to $S_{\text{BH}}$, early-Universe GW \\
\textbf{Tier 3} (strong variants) & $\Delta > 10^{-1}$ & None (Planck-scale) & Spectral dimension, graviton scattering \\
\bottomrule
\end{tabular}
\end{table}

\begin{remark}
\label{rem:tier_tension}
The tension is between Tiers 1 and 2: Tier-1 observables are accessible but weakly distinguishing; Tier-2 observables are more distinguishing but barely accessible. The Quantum Gravity Audit Program is fundamentally about navigating this tension---finding the sweet spot where $\Delta$ is large enough and $M$ is achievable.
\end{remark}

\subsection{The Yajie Consensus at Planck Scale}

The \Yajie{} algorithm for noise detection in data~\cite{scx_theory} has a natural interpretation in the QG audit:

\begin{proposition}[Yajie as QG Consensus]
\label{prop:yajie_qg}
\rigorPartial
In the QG audit context, the \Yajie{} clean data detector becomes a \textbf{gauge-invariant consensus detector}:
\begin{equation}
  \Yajie(\{o_m\}_{m=1}^{M}) =
  \begin{cases}
    \text{``Clean''} & \text{if } \max_{A \in \{\text{ST, LQG, CDT}\}} \Pbb(\{o_m\} \mid A) > \tau,\\[4pt]
    \text{``Dirty''} & \text{otherwise},
  \end{cases}
\end{equation}
where ``clean'' means: the data are consistent with \emph{all} theories in the gauge orbit (the theories are \CI{}), and ``dirty'' means: the data favor one theory over others (the gauge orbit can be split by experiment). The threshold $\tau$ is determined by $M_{\text{crit}}$ and the desired confidence.
\end{proposition}

\subsection{The Spring Gatekeeper for Theory Evolution}

The \Spring{} self-evolution framework~\cite{scx_spring} models the evolution of the QG theory landscape over time. When new experimental data arrive (e.g., from a next-generation CMB experiment), the gatekeeper evaluates whether the data force a ``seasonal transition'' from one theory to another:

\begin{definition}[QG Seasonal Transition]
\label{def:qg_seasonal}
A \textbf{QG seasonal transition} occurs when new experimental data $\D_{\text{new}}$ of size $M_{\text{new}}$ causes the \Cercis{} score ranking to invert:
\begin{equation}
  S_{\QG}(A, B; \Omega_{\text{old}}) < S_{\QG}(B, A; \Omega_{\text{old}})
  \quad \text{but} \quad
  S_{\QG}(A \cup \D_{\text{new}}, B \cup \D_{\text{new}}; \Omega_{\text{exp}}) > \cdots
\end{equation}
If no finite $M_{\text{new}}$ can cause such an inversion, then $A$ and $B$ are in \textbf{eternal winter}: the \Spring{} gatekeeper will never transition between them, confirming \CI{}.
\end{definition}

\subsection{Minimal Distinguishing Experiment Design}

If $M_{\text{crit}}$ is finite and achievable, the audit program specifies the optimal experimental configuration:

\begin{algorithm}[H]
\caption{Optimal QG Distinguishing Experiment Design}
\label{alg:optimal}
\begin{algorithmic}[1]
\Require Theories $A, B$, observable catalog $\Omega_{\text{exp}}$, budget $M_{\text{budget}}$, target $\delta$
\Ensure Experiment configuration $\{(O_m, N_m)\}_{m=1}^{M}$ or \CI{} verdict
\State Compute $\Delta(O)$ for all $O \in \Omega_{\text{exp}}$
\State Sort observables: $\Delta(O_1) \geq \Delta(O_2) \geq \cdots$
\State $\bar{\Delta} \gets 0$, $M \gets 0$
\While{$M < M_{\text{budget}}$ \textbf{and} $\bar{\Delta}^2 < \ln(1/\delta) / (2M)$}
    \State $M \gets M + 1$
    \State $\bar{\Delta} \gets \frac{1}{M} \sum_{m=1}^{M} \Delta(O_m)$
\EndWhile
\If{$M > M_{\text{budget}}$}
    \State \Return \CI{} verdict \Comment{Cannot distinguish within budget}
\Else
    \State \Return $\{(O_m, 1)\}_{m=1}^{M}$ \Comment{One probe per top-$M$ observable}
\EndIf
\end{algorithmic}
\end{algorithm}

\begin{remark}
\label{rem:budget_agnostic}
Algorithm~\ref{alg:optimal} is budget-agnostic: it gives the minimal set of observables needed regardless of cost. In practice, cost constraints may force suboptimal observable choices, increasing the effective $M_{\text{crit}}$.
\end{remark}


% ================================================================
\section{Mathematical Infrastructure: The Audit Hilbert Space}
\label{sec:math_infra}
% ================================================================

\subsection{States, Observables, and Gauge Orbits}

We formalize the QG audit as a quantum statistical decision problem on an abstract Hilbert space that encodes the gauge structure.

\begin{definition}[Audit Hilbert Space]
\label{def:audit_hilbert}
The \textbf{audit Hilbert space} for quantum gravity is:
\begin{equation}
  \Hilb_{\text{audit}} = \bigoplus_{[T] \in \T_{\QG}/\sim_{\G}} \Hilb_{[T]},
\end{equation}
where $\Hilb_{[T]}$ is the observable Hilbert space associated with the gauge equivalence class $[T]$. The direct sum runs over physically distinct equivalence classes.
\end{definition}

\begin{definition}[Observable Algebra on the Audit Space]
\label{def:obs_algebra}
The observable algebra $\A_{\text{audit}} \subset \B(\Hilb_{\text{audit}})$ consists of all bounded operators that are \textbf{gauge-invariant}:
\begin{equation}
  \A_{\text{audit}} = \{O \in \B(\Hilb_{\text{audit}}) : U_g O U_g^{-1} = O,\; \forall g \in \G_{\QG}\},
\end{equation}
where $U_g$ is the unitary (or isometric) representation of the gauge transformation $g$ on $\Hilb_{\text{audit}}$.
\end{definition}

\begin{proposition}[Gauge-Invariant Observables Are Exactly the Experimental Ones]
\label{prop:gauge_inv_experimental}
\rigorPartial
$\A_{\text{audit}} = \A^{\text{(exp)}}$. That is, every gauge-invariant observable is experimentally realizable (in principle), and every experimentally realizable observable is gauge-invariant.
\end{proposition}

\begin{proof}[Proof Sketch]
If $O$ is gauge-invariant, its expectation value depends only on the equivalence class $[T]$, not on the representative. This means $O$ can be measured without knowing which specific QG theory is correct---it is a ``theory-independent'' observable. All experimentally realizable observables have this property, since no experiment has resolved Planck-scale physics. Conversely, if $O$ is experimentally realizable, it must be finite-$M$ resolvable, hence its expectation value is invariant under gauge transformations that leave finite-$M$ distributions invariant.
\end{proof}

\subsection{Information-Theoretic Bounds}

\begin{theorem}[Quantum Fidelity Bound for QG Distinguishability]
\label{thm:fidelity_bound}
\rigorFull
Let $\rho_A^{(M)}$ and $\rho_B^{(M)}$ be the density operators on $\Hilb_{\text{audit}}^{\otimes M}$ representing the predictions of theories $A$ and $B$ for $M$ independent probes. Let $F(\rho, \sigma) = \|\sqrt{\rho}\sqrt{\sigma}\|_1$ be the quantum fidelity. Then the minimum error probability for distinguishing $A$ and $B$ from $M$ probes is bounded below by:
\begin{equation}
  \label{eq:quantum_bound}
  P_{\text{err}}^{(M)} \geq \frac{1}{2}\left(1 - \sqrt{1 - F(\rho_A^{(M)}, \rho_B^{(M)})^2}\right).
\end{equation}
For product states $\rho_A^{(M)} = \rho_A^{\otimes M}$, we have $F(\rho_A^{(M)}, \rho_B^{(M)}) = F(\rho_A, \rho_B)^M$. Hence:
\begin{equation}
  P_{\text{err}}^{(M)} \geq \frac{1}{2}\left(1 - \sqrt{1 - F_0^{2M}}\right),
\end{equation}
where $F_0 = F(\rho_A, \rho_B)$ is the single-probe fidelity. If $F_0 = 1$ (the states are identical), $P_{\text{err}}^{(M)} = 1/2$ for all $M$---the theories are \CI{} in the quantum sense.
\end{theorem}

\begin{proof}
This is the quantum Chernoff bound~\cite{audenaert2007quantum} applied to the two-state discrimination problem. The Helstrom bound gives the minimum error probability $P_{\text{err}} = \frac{1}{2}(1 - \frac{1}{2}\|\rho_A^{(M)} - \rho_B^{(M)}\|_1)$. The fidelity bounds the trace distance: $\frac{1}{2}\|\rho - \sigma\|_1 \leq \sqrt{1 - F(\rho, \sigma)^2}$. Substituting yields the result. For product states, $F(\rho^{\otimes M}, \sigma^{\otimes M}) = F(\rho, \sigma)^M$ by multiplicativity of fidelity under tensor products.
\end{proof}

\begin{corollary}[Quantum $M_{\text{crit}}$]
\label{cor:quantum_mcrit}
For quantum states with fidelity $F_0$, achieving $P_{\text{err}} \leq \delta$ requires:
\begin{equation}
  M \geq M_{\text{crit}}^{\text{(quantum)}} = \left\lceil
  \frac{\ln(1 - (1-2\delta)^2)}{2\ln F_0}
  \right\rceil.
\end{equation}
If $F_0 = 1$, then $M_{\text{crit}}^{\text{(quantum)}} = \infty$, confirming \CI{} from the quantum information perspective.
\end{corollary}


% ================================================================
\section{The Landscape: Implications and the Multiverse}
\label{sec:landscape}
% ================================================================

\subsection{The String Landscape as a Gauge Orbit}

String theory's $\sim 10^{500}$ flux vacua~\cite{bousso2007landscape} are often presented as a ``landscape'' of distinct physical theories. In the gauge-equivalence framework, they are \textbf{points on the same gauge orbit}:

\begin{proposition}[Landscape = Gauge Orbit]
\label{prop:landscape_gauge}
\rigorSketch
All $10^{500}$ flux vacua of string theory belong to the same gauge equivalence class---they are different mathematical representatives of the same physical theory, distinguished only by Planck-scale parameters that leave the observable sector (nearly) invariant.
\end{proposition}

\begin{proof}[Proof Sketch]
The flux vacua differ by the values of stabilized moduli (the compactification manifold's shape and the flux integers). These moduli determine the low-energy effective field theory: the gauge group, the particle spectrum, and the coupling constants. However, the subset of moduli that affect observable physics is tiny: the cosmological constant $\Lambda$, the fine-structure constant $\alpha$, and a handful of Standard Model parameters. The vast majority of the $10^{500}$ vacua differ only in Planck-scale details that are unobservable. Therefore, the effective observable sector $\Obs(T)$ is (nearly) constant across the landscape---the vacua are gauge-equivalent.
\end{proof}

\honestcrit{This is a significant reinterpretation: the landscape problem (``why this vacuum?'') is not a selection problem among physically distinct theories, but a gauge-fixing problem: which mathematical representative we choose to describe the unique observable physics. The anthropic principle becomes a gauge choice.}

\subsection{The Multiverse as Gauge Redundancy}

If the multiverse exists (as suggested by eternal inflation and the string landscape), each ``universe'' with different low-energy parameters is a different point on the \emph{same} gauge orbit. The multiverse is not a collection of physically distinct universes, but a collection of gauge-equivalent descriptions of the same underlying physics, projected differently onto the observable sector in each ``branch.''

\begin{conjecture}[Multiverse CI]
\label{conj:multiverse_ci}
All ``universes'' in the eternally inflating multiverse (or the string landscape) belong to a single gauge equivalence class $[T_{\text{QG}}] \in \T_{\QG}/\sim_{\G}$. The apparent variation in low-energy parameters is an artifact of projecting the same gauge-invariant state onto different local observable algebras.
\end{conjecture}

\begin{remark}
\label{rem:multiverse_ci}
If Conjecture~\ref{conj:multiverse_ci} holds, then the multiverse is \CI{}-inseparable from a single universe from the perspective of any finite observer. The question ``why this vacuum?'' is ill-posed: all vacua are gauge-equivalent, and the perception of being in a particular vacuum is a gauge choice made by the observer's finite experimental horizon.
\end{remark}


% ================================================================
\section{Connections to Existing Frameworks}
\label{sec:connections}
% ================================================================

\subsection{Gauge-Gravity Duality and the Audit Limit}

\begin{remark}[The Audit Limit as a Holographic Screen]
\label{rem:audit_screen}
The finite-$M$ audit limit ($M \leq M_{\text{universe}}$) functions as a \textbf{holographic screen} for quantum gravity experiments: it bounds the amount of information about the Planck-scale microstructure that can reach a finite observer. This is reminiscent of the covariant entropy bound~\cite{bousso2002holographic}: the information accessible on a light-sheet is bounded by the area of its boundary in Planck units. The audit limit refines this: even if information \emph{exists} in the full theory, it may not be \emph{experimentally resolvable} with finitely many probes.
\end{remark}

\subsection{Quantum Reference Frames}

The QG gauge group $\G_{\QG}$ is a symmetry group in the sense of quantum reference frames~\cite{giacomini2019quantum}: choosing a representative theory (string, LQG, or CDT) is analogous to choosing a quantum reference frame for describing the same physics. Just as a quantum particle's state looks different in different quantum reference frames but represents the same physical situation, quantum gravity looks different in string, LQG, or CDT variables but represents the same physical theory on the observable sector.

\subsection{Emergent Spacetime and the Gauge Orbit}

If spacetime is emergent (as suggested by AdS/CFT, entanglement gravity, and some formulations of LQG and CDT), then the gauge orbit acquires a deeper meaning: the different QG theories are different descriptions of \emph{how} spacetime emerges, not different predictions about \emph{what} emerges. The emergent spacetime is the gauge-invariant; the microphysical description (strings, spin networks, simplices) is the gauge-variant.

\begin{proposition}[Emergence $\implies$ Gauge Equivalence]
\label{prop:emergence_gauge}
\rigorSketch
If spacetime emerges from a more fundamental description, then all candidate emergence mechanisms that produce the same low-energy effective geometry are gauge-equivalent. The choice of microphysical ``building blocks'' (strings, spin networks, simplices) is a gauge choice.
\end{proposition}


% ================================================================
\section{Discussion}
\label{sec:discussion}
% ================================================================

\subsection{Summary}

We have introduced the \textbf{Compactness-Inseparability (CI) criterion} for quantum gravity theories and argued that string theory, loop quantum gravity, and causal dynamical triangulations are either exactly \CI{} (gauge-equivalent on the observable sector) or \CI{}-borderline (distinguishable only with a probe budget exceeding practical cosmological limits).

The key results are:
\begin{enumerate}[leftmargin=*, label=(\arabic*)]
  \item The space of QG theories admits a gauge group $\G_{\QG}$ whose action preserves all finite-$M$ observable distributions (Theorem~\ref{thm:gauge_equivalence}).
  \item $S_{\text{BH}} = A/4\ell_P^2$ is a gauge invariant; cosmological predictions are gauge-variant but experimentally borderline (Theorem~\ref{thm:cercis_qg}, Proposition~\ref{prop:cosmo_variant}).
  \item $M_{\text{crit}}$ has a closed form $M_{\text{crit}} = \lceil \ln(1/\delta) / (2\bar{\Delta}^2) \rceil$ (Theorem~\ref{thm:m_crit}).
  \item AdS/CFT is the $M \to \infty$ perfect audit with zero experimental resolvability (Theorem~\ref{thm:adscft_audit}).
  \item The \CI{} criterion is equivalent to logical compactness: two theories are \CI{} iff the union of their prediction sets is finitely satisfiable (Theorem~\ref{thm:ci_logic}).
\end{enumerate}

\subsection{What This Paper Does \emph{Not} Claim}

It is important to state clearly what this paper does \emph{not} claim:

\begin{itemize}[leftmargin=*]
  \item \textbf{We do not claim to have proven} that ST, LQG, and CDT are gauge-equivalent. Conjecture~\ref{conj:unification} is exactly that---a conjecture. What we have provided is the \emph{framework} for testing it.
  \item \textbf{We do not claim} that quantum gravity is unobservable in principle. The \CI{} criterion has a sharp boundary: if $\bar{\Delta} > 0$, there exists a finite $M_{\text{crit}}$ (however large) that can distinguish the theories. The question is whether $M_{\text{crit}} \leq M_{\text{universe}}$.
  \item \textbf{We do not claim} that all QG theories are equivalent. The gauge equivalence class is defined by the observable sector; theories with different observable predictions belong to different classes. The claim is that ST, LQG, and CDT are in the \emph{same} class.
\end{itemize}

\subsection{Honest Limitations}

\honestcrit{Several serious limitations temper our results:}

\begin{enumerate}[leftmargin=*, label=(\arabic*)]
  \item \textbf{No explicit gauge transformations.} We have not constructed the maps $(\phi, \psi)$ that connect ST to LQG or CDT. This is a formidable mathematical challenge that may require a non-perturbative formulation of all three theories---something that currently does not exist.
  \item \textbf{The $\bar{\Delta}$ estimates are order-of-magnitude.} The per-observable distinguishabilities in Table~\ref{tab:mcrit_estimates} are rough estimates. A rigorous computation would require full control over the theory parameter spaces, including measure factors over the string landscape and quantization ambiguities in LQG.
  \item \textbf{$M_{\text{universe}}$ is not a sharp bound.} The cosmological probe budget is a heuristic, not a theorem. Different experimental modalities (CMB, GW, astrophysical BH, laboratory analogues) are not fully independent, and their correlations reduce the effective $M$.
  \item \textbf{The logical compactness connection is formal.} While Theorem~\ref{thm:ci_logic} establishes a precise correspondence between \CI{} and first-order compactness, the language $\Lang_{\QG}$ is a simplified model of real experimental observables. The actual ``language'' of quantum gravity phenomenology is much richer.
  \item \textbf{The gauge group $\G_{\QG}$ may not be a group.} In general, the gauge transformations between different QG theories may form a groupoid (as in Definition~\ref{def:gauge_group}) rather than a group, since composition may not always be defined. This complicates the orbit space structure.
\end{enumerate}

\subsection{Future Directions}

\begin{enumerate}[leftmargin=*, label=(\arabic*)]
  \item \textbf{Construct explicit gauge transformations.} The first priority is to find concrete maps between ST, LQG, and CDT Hilbert spaces that preserve observable expectation values. This may be possible in simplified models (e.g., 3D gravity, where all three frameworks are better understood).
  \item \textbf{Precision cosmology as QG audit.} Next-generation CMB experiments (Simons Observatory, CMB-S4) and GW observatories (LISA, Einstein Telescope) will shrink the error bars on $n_s$, $r$, and the stochastic GW background. Each improvement reduces $M_{\text{crit}}$ and pushes the \CI{} boundary.
  \item \textbf{Laboratory analogues.} Condensed matter systems (cold atoms, superconducting circuits) can simulate aspects of quantum gravity~\cite{bravetti2015analogue}. These provide $M \gg 1$ in a controlled setting, albeit for analogue rather than real gravity.
  \item \textbf{Quantum information constraints.} The quantum fidelity bound (Theorem~\ref{thm:fidelity_bound}) can be sharpened using more sophisticated quantum hypothesis testing techniques, potentially yielding a tighter $M_{\text{crit}}^{\text{(quantum)}}$.
  \item \textbf{The \Spring{} gatekeeper for theory space.} The self-evolution framework of \Spring{} can be applied to the QG theory landscape itself: as experimental data accumulate, the gatekeeper evolves, and theories are selected or discarded based on \Cercis{} score trajectories.
\end{enumerate}

\subsection{Philosophical Coda}

The \CI{} criterion reframes the quantum gravity problem in a way that is simultaneously humbling and clarifying. If ST, LQG, and CDT are gauge-equivalent, then the decades-long debate about which theory is ``correct'' is a debate about gauge choices---a category error. The real question is not ``which theory is right?'' but ``how many independent experimental probes are needed to see through the gauge to the invariant physics?''

If the answer is $M_{\text{crit}} > M_{\text{universe}}$, then quantum gravity is \CI{} to us---operationally indistinguishable by any experiment we can ever perform within this universe. This would not make quantum gravity any less real; it would make it a gauge theory in the deepest sense, where the gauge freedom extends all the way to the horizon of experimental accessibility.

The \SCX{} framework provides the mathematical language for this reframing: the \Cercis{} score between theories, the \Yajie{} consensus audit, and the \Spring{} gatekeeper for theory evolution. Quantum gravity is not a problem to be solved by choosing the right theory---it is an audit to be conducted by accumulating the right probes.


% ================================================================
% Appendix: Glossary of Notation
% ================================================================
\section*{Appendix A: Glossary of Notation}
\label{app:notation}

\begin{table}[H]
\centering
\begin{tabular}{@{}ll@{}}
\toprule
\textbf{Symbol} & \textbf{Meaning} \\
\midrule
$T = (\Hilb_T, \A_T, \rho_T)$ & Quantum gravity theory \\
$\Obs(T)$ & Observable sector of theory $T$ \\
$\G_{\QG}$ & QG gauge groupoid \\
$\sim_{\G}$ & Gauge equivalence relation \\
$\T_{\QG} / \sim_{\G}$ & Space of gauge equivalence classes \\
$S_{\QG}(A, B; \Omega)$ & QG \Cercis{} score between theories $A$ and $B$ \\
$Q_{\text{agree}}$ & Agreement (quality) component of \Cercis{} score \\
$N_{\text{dist}}$ & Distinguishability (novelty) component \\
$\eta(M)$ & Resolution schedule \\
$\Delta(O) = \TV(\Pbb_A, \Pbb_B)$ & Per-observable distinguishability \\
$\bar{\Delta}$ & Average distinguishability over observable ensemble \\
$M_{\text{crit}}(A, B; \delta)$ & Minimum probes for confidence $1-\delta$ \\
$M_{\text{universe}}$ & Cosmological probe budget \\
\CI{} & Compactness-Inseparability \\
$d_{\G}(A, B)$ & Gauge orbit diameter \\
$F(\rho, \sigma)$ & Quantum fidelity \\
$\Hilb_{\text{audit}}$ & Audit Hilbert space \\
$\A_{\text{audit}}$ & Gauge-invariant observable algebra \\
\bottomrule
\end{tabular}
\end{table}


% ================================================================
% Appendix: Numerical Verification
% ================================================================
\section*{Appendix B: Numerical Verification of $M_{\text{crit}}$}
\label{app:numerical}

The Python script \texttt{verify\_qg\_audit.py} accompanying this paper provides numerical verification of the key quantitative claims:

\begin{enumerate}[leftmargin=*, label=(\arabic*)]
  \item Computation of $M_{\text{crit}}$ from $\bar{\Delta}$ and $\delta$ using Eq.~\eqref{eq:m_crit}.
  \item Monte Carlo simulation of the hypothesis testing procedure for two theories with specified $\bar{\Delta}$.
  \item Empirical verification that the error probability matches the Hoeffding bound.
  \item \CI{} boundary detection: identifying when $M_{\text{crit}}$ exceeds specified budget.
  \item \Cercis{} score computation for pairwise theory comparisons using Table~\ref{tab:cercis_map} data.
\end{enumerate}

All numerical claims in the paper can be reproduced by running:
\begin{verbatim}
python verify_qg_audit.py
\end{verbatim}
with \texttt{numpy} and \texttt{scipy} installed.


% ================================================================
\begin{thebibliography}{99}

\bibitem{rovelli2004quantum}
C.~Rovelli, \emph{Quantum Gravity}, Cambridge University Press, 2004.

\bibitem{polchinski1998string}
J.~Polchinski, \emph{String Theory}, Cambridge University Press, 1998.

\bibitem{ambjorn2012nonperturbative}
J.~Ambj{\o}rn, A.~G\"orlich, J.~Jurkiewicz, and R.~Loll,
\emph{Nonperturbative quantum gravity}, Phys.~Rep.~519 (2012), 127--210.

\bibitem{green1987superstring}
M.~B.~Green, J.~H.~Schwarz, and E.~Witten,
\emph{Superstring Theory}, Cambridge University Press, 1987.

\bibitem{ashtekar2004background}
A.~Ashtekar and J.~Lewandowski,
\emph{Background independent quantum gravity: A status report},
Class.~Quant.~Grav.~21 (2004), R53.

\bibitem{ambjorn2001nonperturbative}
J.~Ambj{\o}rn, J.~Jurkiewicz, and R.~Loll,
\emph{Nonperturbative Lorentzian quantum gravity, causality and topology change},
Nucl.~Phys.~B~610 (2001), 347--382.

\bibitem{strominger1996microscopic}
A.~Strominger and C.~Vafa,
\emph{Microscopic origin of the Bekenstein--Hawking entropy},
Phys.~Lett.~B~379 (1996), 99--104.

\bibitem{ashtekar2006quantum}
A.~Ashtekar, T.~Pawlowski, and P.~Singh,
\emph{Quantum nature of the big bang},
Phys.~Rev.~Lett.~96 (2006), 141301.

\bibitem{ambjorn2005spectral}
J.~Ambj{\o}rn, J.~Jurkiewicz, and R.~Loll,
\emph{Spectral dimension of the universe},
Phys.~Rev.~Lett.~95 (2005), 171301.

\bibitem{maldacena1999large}
J.~Maldacena,
\emph{The large $N$ limit of superconformal field theories and supergravity},
Int.~J.~Theor.~Phys.~38 (1999), 1113--1133.

\bibitem{hawking1976black}
S.~W.~Hawking,
\emph{Black holes and thermodynamics},
Phys.~Rev.~D~13 (1976), 191--197.

\bibitem{strominger2001ds}
A.~Strominger,
\emph{The dS/CFT correspondence},
JHEP~10 (2001), 034.

\bibitem{baumann2015inflation}
D.~Baumann and L.~McAllister,
\emph{Inflation and String Theory}, Cambridge University Press, 2015.

\bibitem{bousso2007landscape}
R.~Bousso and J.~Polchinski,
\emph{The string theory landscape},
Sci.~Am.~291 (2004), 60--69.

\bibitem{bousso2002holographic}
R.~Bousso,
\emph{The holographic principle},
Rev.~Mod.~Phys.~74 (2002), 825--874.

\bibitem{giacomini2019quantum}
F.~Giacomini, E.~Castro-Ruiz, and \v{C}.~Brukner,
\emph{Quantum mechanics and the covariance of physical laws in quantum reference frames},
Nature~Commun.~10 (2019), 494.

\bibitem{audenaert2007quantum}
K.~M.~R.~Audenaert et al.,
\emph{Discriminating states: The quantum Chernoff bound},
Phys.~Rev.~Lett.~98 (2007), 160501.

\bibitem{bravetti2015analogue}
C.~Barcel\'o, S.~Liberati, and M.~Visser,
\emph{Analogue gravity},
Living~Rev.~Rel.~14 (2011), 3.

\bibitem{scx_theory}
SCX,
\emph{State-Conditioned Expertise: A Complete Theory of Label Noise Detection via Multi-Expert Consistency},
SCX Technical Report, 2026.

\bibitem{scx_compactness}
SCX,
\emph{Compactness and SCX: From Finite to Infinite Audit Transfer},
SCX Technical Report, 2026.

\bibitem{scx_spring}
SCX,
\emph{Spring: A Self-Evolving Gatekeeper with Provable Convergence},
SCX Technical Report, 2026.

\end{thebibliography}

\end{document}
